\documentclass[11pt,a4paper]{article}
\usepackage[T1]{fontenc}
\usepackage[utf8]{inputenc}
\usepackage{lmodern,microtype}
\usepackage[margin=.85in]{geometry}
\usepackage{amsmath,amssymb,amsthm,mathtools,aliascnt}
\usepackage{booktabs,tabularx,array,enumitem,needspace,float}
\usepackage[section]{placeins}
\usepackage{xcolor,tikz,mdframed}
\usetikzlibrary{arrows.meta,positioning,calc}
\usepackage[hidelinks]{hyperref}
\usepackage[nameinlink,noabbrev]{cleveref}
\usepackage{fancyhdr}
\hypersetup{pdftitle={Sync HotStuff RR v3.3},pdfsubject={Proposal timers and terminal-QC handoff}}
\setlength{\headheight}{14pt}
\pagestyle{fancy}\fancyhf{}
\lhead{\small Sync HotStuff RR v3.3}\rhead{\small L1-scheduled leaders}
\cfoot{\thepage}
\setlength{\parindent}{0pt}\setlength{\parskip}{4pt}
\setlength{\emergencystretch}{2em}
\setlist{nosep,leftmargin=*}
\newcolumntype{Y}{>{\raggedright\arraybackslash}X}
\definecolor{rrRed}{HTML}{A51C30}
\definecolor{rrBlue}{HTML}{2563EB}
\definecolor{rrTeal}{HTML}{0F766E}
\newcommand{\RR}[1]{{\color{rrRed}#1}}
\newcommand{\QC}{\mathcal C}
\newcommand{\Hash}{\mathsf H}
\newcommand{\rank}{\operatorname{rank}}
\newcommand{\tagg}[1]{\ensuremath{\mathsf{#1}}}
\newtheorem{definition}{Definition}[section]
\newaliascnt{assumption}{definition}
\newtheorem{assumption}[assumption]{Assumption}
\aliascntresetthe{assumption}
\newaliascnt{lemma}{definition}
\newtheorem{lemma}[lemma]{Lemma}
\aliascntresetthe{lemma}
\newaliascnt{theorem}{definition}
\newtheorem{theorem}[theorem]{Theorem}
\aliascntresetthe{theorem}
\newaliascnt{proposition}{definition}
\newtheorem{proposition}[proposition]{Proposition}
\aliascntresetthe{proposition}
\crefname{assumption}{assumption}{assumptions}
\Crefname{assumption}{Assumption}{Assumptions}
\crefname{lemma}{lemma}{lemmas}
\Crefname{lemma}{Lemma}{Lemmas}
\crefname{theorem}{theorem}{theorems}
\Crefname{theorem}{Theorem}{Theorems}
\newcounter{protocol}
\newenvironment{protocol}[2]{\refstepcounter{protocol}\label{#1}
\begin{mdframed}[nobreak=true,linewidth=.6pt,skipabove=8pt,skipbelow=8pt,
innertopmargin=7pt,innerbottommargin=7pt]\small
\textbf{Protocol \theprotocol. #2}\par\smallskip}{\end{mdframed}}
\crefname{protocol}{Protocol}{Protocols}
\Crefname{protocol}{Protocol}{Protocols}
\title{Sync HotStuff RR v3.3\\\large Proposal Timers and Terminal-QC Handoff}
\author{}\date{12 October 2026}
\begin{document}
\maketitle
\begin{abstract}
An L1 authority assigns the leader of each generation and term.
A term certifies $K\ge1$ new DATA blocks, called fresh blocks. Its
terminal certificate admits the next term without waiting for commitment.
Replicas forward that certificate to the next leader as a backup to
ordinary vote delivery. Each DATA proposal that a replica votes for
starts one local $2\Delta$ commit timer at that replica. This includes a
previously unseen proposal from an earlier term of the same generation.
No certificate starts a separate timer. A complete quit proof freezes the
generation and invokes Sync HotStuff's lock-based recovery into $(g+1,0)$.
We prove safety under authenticated synchrony and eventual liveness under
explicit L1 availability, permanent equivocation exclusion, and fairness
assumptions. We also give counterexamples to earlier-term vote skipping,
shortened progress deadlines, and unconditional liveness without exclusion.
\end{abstract}

\section{Purpose and replication guarantees}
\label{sec:scope}
This specification revises the standard-synchronous Sync HotStuff
protocol~\cite{sync}. It separates ordinary leader rotation from recovery:
\[
 (g,t)\xrightarrow{\text{terminal DATA QC}}(g,t+1),
 \qquad
 g\xrightarrow{\text{complete quit proof}}(g+1,0).
\]
Ordinary rotation preserves commit timers. Recovery cancels all unfinished
timers in the departing generation. Compared with v3, v3.1 removes
QC-start timers and vote skipping, gives explicit parent-selection rules,
and adds terminal-QC forwarding to the next leader.
 The progress checks of \cref{alg:recovery} start at
the term-entry time recorded by \cref{alg:steady}.

This version keeps the intended protocol of v3.1. It moves rules from
the prose into the protocol boxes, defines terms that v3.1 used without
definition, and fixes notation. Where v3.1 allowed two readings, this
version states the reading that the proofs use.

\begin{definition}[State machine replication]
Blocks order commands for a deterministic state
machine~\cite{schneider}. Committing a
block commits its uncommitted ancestors in order. \emph{Safety} requires
correct replicas' committed logs to be prefix-comparable and never revoked.
\emph{Validity} requires application-valid committed blocks.
\emph{Request liveness} requires every correct-client request that remains
admissible until ordered to be committed eventually by every correct replica.
\end{definition}
We first prove unbounded committed-height progress, then derive request
liveness using an explicit input-fairness assumption. Normally completed
terms contain $K$ fresh blocks. A faulty term can produce none. Then BLAME freezes its generation, and recovery starts the
next generation. Terms are defined by certified block counts, not
fixed physical durations. Replicas use a global clock but may enter a term
at different times.

\section{Model and service assumptions}
\label{sec:model}
There are $n=2f+1$ fixed voting members, at most $f$
Byzantine. The Byzantine set is fixed from the start.
Initialization uses a common genesis and initial authority; correct
replicas enter $(0,0)$ within $\Delta$ of one another.

\Needspace{9\baselineskip}
\begin{assumption}[Authenticated synchronous communication and global clock]
\label{ass:network}
\leavevmode\par
\begin{itemize}
    \item Messages use authenticated communication channels.
    \item The actual delay of a correct-to-correct message is $\delta\le\Delta$;
    $\delta$ may differ across messages.
    \item All replicas share a global clock.
\end{itemize}
\end{assumption}
The delay bound includes delivery and processing of complete validation
dependencies. Each replica records its own entry and voting events on
the global clock.

\begin{assumption}[Finalized authority]
\label{ass:authority}
L1 is a public bulletin board. For each generation $g$, it supplies one
immutable authority $\tagg{AUTH}_g$ specifying a unique leader $L(g,t)$
from the fixed committee for every term $t$. At the finite global time
$d(g)$, that authority is finalized and simultaneously available to all
correct replicas, including frozen ones. L1 progresses independently
of L2 commitment.
\end{assumption}
An authority may be finalized in advance. The L1 epoch and recovery
generation need not coincide.

\begin{assumption}[Accountability]
\label{ass:exclusion}
Every admissible leader-equivocation proof revealed to a correct replica
is eventually submitted to L1 and processed. For each exposed identity $p$,
let $g_{\mathrm{bad}}$ be the generation identified by the first such proof
revealed to a correct replica. This is the generation of the proved
equivocation, not necessarily the generation current at revelation.
There is a finite removal generation $G_p>g_{\mathrm{bad}}$ such that
\[
 L(g,t)\ne p\quad\text{for every }g\ge G_p\text{ and every }t\ge0.
\]
Exclusion is permanent. The finite Byzantine identity set cannot be
replenished by new corruptions, joins, or replacement identities.
\end{assumption}
Both $g_{\mathrm{bad}}$ and $G_p$ are generation indices. No fixed bound
on $G_p-g_{\mathrm{bad}}$ is assumed. Exclusion takes effect in future
authorities without rewriting a finalized assignment.
This is an enforced eligibility property, not an economic claim that a
deposit deters faults. Timeout BLAME alone is not equivocation evidence.
Excluded proposers may remain faulty voters within the same $f$ bound.
Changing voting membership or quorum weights needs a separate
cross-configuration argument.

\begin{assumption}[Schedule and input fairness]
\label{ass:fairness}
For every $g_0$, some finite $g\ge g_0$ has correct $L(g,0)$.
For every fixed $g$ and $t_0$, some finite $t\ge t_0$ has correct $L(g,t)$.
Correct leaders can propose valid empty blocks.

Correct-client requests eventually reach all correct replicas. They have
a common admission order with finitely many predecessors per request, and
each request fits in a batch. Correct leaders include the oldest pending
admissible requests that fit.  A request stays pending until
it is committed. A leader skips a pending request that is already in the
new block's uncommitted ancestry. Correct replicas retain and serve
validation data.
\end{assumption}
Safety uses \cref{ass:network,ass:authority}. Liveness additionally
uses \cref{ass:exclusion,ass:fairness}.

\section{Blocks, certificates, and message handling}
\label{sec:definitions}
\paragraph{Terms and positions.}
A term is identified by $(g,t)$. Terms are ordered
lexicographically. Generation $g$ has an opening block $A_g$; $A_0$
is genesis. Each term has $K$ fresh positions $j\in\{0,\ldots,K-1\}$.
A fresh block is a DATA block at a fresh position. The opening block is not
fresh. Fresh position $j$ of term $(g,t)$ has height
\[
 k=h(A_g)+tK+j+1.
\]
Its block is $B_k=(x_k,\Hash(B_{\mathrm{parent}}))$, with
$h(B_{\mathrm{parent}})=k-1$, where $x_k$ is a batch of requests.
The term's terminal block has $j=K-1$.
The initializer of generation $g$ is $L(g,0)$. The successor of term
$(g,t)$ is $L(g,t+1)$. A replica's current term is the latest term it has
entered.

\paragraph{Signed context and certificates.}
Every signature binds the protocol, the voting configuration, the L1
authority, the message type, and the term $(g,t)$. DATA and ANCHOR
signatures also bind the generation opening and the block's height and
hash. DATA signatures also bind position $j$. BLAME, STATUS, and NEWVIEW
signatures bind no opening, because none may exist yet.
 We omit these fixed
context fields from message displays.

A DATA QC $\QC_{g,t}(B_k)$ contains $f+1$ matching votes from distinct
members. An ANCHOR QC $\QC_{g,0}(A_g)$ contains $f+1$ matching ANCHOR
votes from distinct members. It certifies, in generation $g$, the opening
block that the NEWVIEW's older QC already certified.
ANCHOR votes are separately tagged. A correct replica casts at most one DATA vote per $(g,t,j)$ and one
ANCHOR vote per generation.

Certificate rank is $(g,t,k)$, where $(g,t)$ is the term in which the
certificate was formed and $k$ is the height of the certified block. The
ANCHOR QC of generation $g$ has rank $(g,0,h(A_g))$. This is below the rank
$(g,0,h(A_g)+1)$ of the generation's first DATA QC.
Equal-rank choices use a fixed total tie-break on certified content,
independently of signer subsets or encodings. Genesis has a synthetic
certificate below all ordinary ranks. An opening's identity is
$(g,h(A_g),\Hash(A_g))$.

\paragraph{Admission and complete validation.}
Generation $0$ is admitted at initialization. Generation $g>0$ is
admitted at replica $r$ once $r$ holds $\tagg{AUTH}_g$ and a quit proof for
generation $g-1$. That quit proof is the \emph{entry proof} of $g$. Term
$(g,0)$ is admitted together with its generation, even before its leader
sends NEWVIEW. Term $(g,t)$ with $t>0$ is admitted once $r$ also holds the
terminal QCs of $(g,0),\ldots,(g,t-1)$. A message is \emph{admissible} if
its term is admitted and its signatures are valid.


For $g>0$, the opening is a signed NEWVIEW from $L(g,0)$, justified by
a valid QC from a strictly older generation and the entry proof for
$g$. Its ANCHOR QC admits fresh DATA proposals. Genesis needs no ANCHOR
exchange.

Validation checks signatures, authority, the opening and term admission,
position and height, the parent rule in \cref{tab:parent}, parent QCs,
hash links, and application validity throughout the required ancestry.
Leader authorization is checked against the public L1 authority.
Every relevant ancestor includes its signed proposal and block data;
a chain of certificate hashes alone is insufficient. Dependencies can
come from the message or a validated local cache. Correct senders provide
complete validation bundles. References are allowed only when the receiver
can resolve them within the delay bound $\Delta$ of
\cref{ass:network}. Bundles list each dependency once rather than
recursively copying nested bundles.

\paragraph{Evidence takes priority.}
Before validating an enclosing message or discarding a duplicate,
inspect all signed material in it for a quit proof (defined
after \cref{alg:steady}). A proof needs its own valid signatures and
admission witnesses. A valid proof counts even when the message that
carries it is invalid or incomplete. Two conflicting messages form a proof
only when both are validly signed and admissible.
 Compare incoming
material with all retained material, not just with other
objects inside the message.

\paragraph{One proposal-processing rule.}
Every previously unvoted, admitted current-generation DATA proposal invokes
the Vote step, including proposals in attached histories, proposals from
earlier terms, and proposals whose blocks have already been committed as
ancestors.  An earlier term does not make a proposal
ineligible. A proposal's identity is $(g,t,j,\Hash(B_k))$. Two
proposals that differ only in the signer set or encoding of
$Q_{\mathrm{parent}}$ are the same proposal. Deduplicate proposals
by this identity after evidence inspection. Old-generation DATA proposals
are retained but do not receive new votes or timers. Future-generation
messages wait for authenticated entry.

\paragraph{Active and frozen generations.}
New proposals, DATA/ANCHOR votes, and ordinary term transitions require
the current generation to be \tagg{Unfrozen}. A frozen generation still
accepts evidence, old votes, and QCs, observes the L1 bulletin board, and
executes recovery. Freezing never undoes a commit.

\begin{table}[htbp]\centering\small
\caption{Parent selection for the next fresh position $j$ of term $(g,t)$.}
\label{tab:parent}
\begin{tabularx}{\linewidth}{@{}l Y Y@{}}
\toprule
Condition & $B_{\mathrm{parent}}$ & $Q_{\mathrm{parent}}$\\
\midrule
$j=0,t=0$ & Opening block $A_g$ & Its ANCHOR QC in $(g,0)$\\
$j=0,t>0$ & Position $K-1$ of $(g,t-1)$ & Its DATA QC in $(g,t-1)$\\
$1\le j<K$ & Position $j-1$ of $(g,t)$ & Its DATA QC in $(g,t)$\\
\bottomrule
\end{tabularx}
\smallskip
For $(g,t,j)=(0,0,0)$, use genesis and its synthetic certificate instead.
\end{table}

\section{Protocols}
\label{sec:protocols}
\RR{Red text} marks extensions relative to standard-synchronous
Sync HotStuff~\cite{sync}. Explanations stay outside the protocol
boxes.

\paragraph{Message-processing order.}
Replica $r$ handles each received message in this order:
\begin{enumerate}[label=(\arabic*)]
\item check all signed material for a quit proof;
\item check completeness and validity;
\item process attached DATA proposals under Vote, parents first;
\item process attached terminal QCs under step~4 of \cref{alg:steady}, in
increasing term order;
\item process the message itself.
\end{enumerate}
When a message and a timer fall at the same time, $r$ handles the message
first.

\begin{protocol}{alg:steady}{Propose, vote, commit, and normal rotation}
Replica $r$ runs these steps. Its current term is $(g,t)$.
\begin{enumerate}[label=\textbf{\arabic*)},itemsep=5pt]
\item \textbf{Propose.} Run this step only if $r=L(g,t)$
\RR{and $g$ is unfrozen}.
For each position $j=0,\ldots,K-1$, in order:
 \begin{itemize}
 \item Wait until $r$ holds $B_{\mathrm{parent}}$ and
       $Q_{\mathrm{parent}}$ for $j$ \RR{(\cref{tab:parent})}.
 \item Build $B_k=(x_k,\Hash(B_{\mathrm{parent}}))$.
       The batch $x_k$ follows \cref{ass:fairness}.
 \item Broadcast
       $P=\langle\tagg{PROPOSE},g,t,j,B_k,Q_{\mathrm{parent}}\rangle_r$
       with its validation bundle.
 \item Process $P$ under Vote.
 \end{itemize}
 Do not wait for the parent to commit. Never sign two different
 blocks for one position. A resend repeats the same $P$.

\item \textbf{Vote.} Run this step on every DATA proposal $P$
of term $(g_P,t_P)$ at position $j_P$. \RR{Here $t_P$ may be earlier
than $t$. The step also applies to proposals inside attached histories.}
 \begin{enumerate}[label=(\alph*)]
 \item \RR{If $P$ or its attached material contains a quit proof for $g$,
       run Quit and stop.}
 \item If $g_P<g$, retain $P$ and stop. If $g_P>g$, buffer $P$
       until $r$ enters generation $g_P$.
 \item \RR{If $g$ is frozen, retain $P$ and stop.}
 \item If a dependency is missing, retain $P$, request the dependency,
       and stop. Rerun this step when it arrives.
 \item If $P$ is invalid, discard it and stop.
 \item If $r$ has already voted for $P$, stop.
       \RR{A duplicate never restarts a timer.}
 \item Otherwise, let $\tau_r(P)$ be the current time. Forward $P$ and
       its bundle to all replicas. Broadcast
       $\langle\tagg{DATA\mbox{-}VOTE},g,t_P,j_P,\Hash(B_k)\rangle_r$.
       Start the commit timer of $P$. It expires at $\tau_r(P)+2\Delta$.
 \end{enumerate}

\item \textbf{Commit.} When the timer of $P$ expires, commit $B_k$ and
its uncommitted ancestors. \RR{Only Quit cancels a timer. Entering a
new term does not cancel it. A QC never starts a timer.}

\item \textbf{Terminal QC and rotation.} \RR{Run this step the first time
$r$ obtains a valid terminal QC for a term $(g,t')$.}
\RR{``Obtains'' means that $r$ assembles the QC from votes or finds it in
any message, including as another proposal's
$Q_{\mathrm{parent}}$.}
 \begin{enumerate}[label=(\alph*)]
 \item \RR{Process each unseen proposal in the QC's history under Vote,
       parents first.}
 \item \RR{If $g$ is now frozen, stop.}
 \item \RR{Send the QC and its validation bundle to $L(g,t'+1)$.}
 \item \RR{If $t\le t'$, enter $(g,t'+1)$ and set $e_r(g,t'+1)$ to the
       current time.} \RR{Cancel the progress checks of all earlier
       terms.}
 \end{enumerate}
 \RR{A QC never moves $r$ to an earlier term. It never changes a recorded
 entry time.}
\end{enumerate}
\end{protocol}

The acceptance time $\tau_r(P)$ is when complete validation and the first
vote occur, not an earlier arrival time of an incomplete packet.
Each replica records its own entry time for each term. Its entry time
$e_r(0,0)$ is its initialization time. For $g>0$, $e_r(g,0)$ is set in
step~3 of \cref{alg:recovery}. For $t>0$, $e_r(g,t)$ is set in step~4 of
\cref{alg:steady}. A term's progress checks therefore start when the term
is admitted locally, not when its leader's first proposal arrives. A
silent new leader cannot delay them.


\paragraph{Complete quit proofs.}
For an admitted generation $g$, a proof consists of either:
\begin{enumerate}[label=(\alph*)]
\item $f+1$ distinct signatures on the same $\langle\tagg{BLAME},g,t\rangle$
for an admitted term, with that term's admission witness;
\item two different admissible DATA blocks signed by $L(g,t)$ at the same
$(g,t,j)$, with their admission histories; or
\item two valid NEWVIEWs signed by $L(g,0)$ selecting different opening
blocks for generation $g$.
\end{enumerate}
Different openings count even when their blocks are comparable, because
they fix different starting heights $h(A_g)$. The height formula then
maps the same position to different heights.

Different QCs for the same opening block are not equivocation. Evidence
from any admitted term can freeze its generation, including a term not yet
locally reached. Old-generation proofs are archival and may be submitted to
L1. Future-generation proofs are buffered until entry. Neither resets an
unrelated recovery.

\begin{protocol}{alg:recovery}{Blame, quit, and recovery}
Replica $r$ runs these steps. Its current term is $(g,t)$.
\begin{enumerate}[label=\textbf{\arabic*)},itemsep=5pt]
\item \textbf{Blame.} For each $i=1,\ldots,K$, run a progress check at
time
\[
 e_r(g,t)+(2i+4)\Delta.
\]
At that time, broadcast $\langle\tagg{BLAME},g,t\rangle_r$ if all of these
hold:
 \begin{itemize}
 \item $r$'s current term is still $(g,t)$, and $g$ is unfrozen;
 \item $r$ has cast DATA votes for fewer than $i$ positions of term
       $(g,t)$ \RR{(ANCHOR votes do not count)};
 \item $r$ has not yet sent BLAME for $(g,t)$.
 \end{itemize}

\item \textbf{Quit.} Run this step the first time $r$ holds a quit proof
for its current generation $g$.
 \begin{itemize}
 \item \RR{Freeze $g$. Stop all progress checks, proposals, votes, and
       term changes in $g$. Cancel every unexpired commit timer of $g$.}
 \item Set $a_r$ to the current time. Broadcast the proof.
 \item Keep all commitments. Keep collecting votes and QCs of $g$.
       They can raise the entry lock chosen in step~3.
 \item \RR{If the proof is of type (b) or (c), submit it to L1. Do
       not wait for L1 to process it.}
 \end{itemize}

\item \textbf{Status and entry.} \RR{At time
$e_r(g+1,0)=\max\{a_r+\Delta,d(g+1)\}$, read $\tagg{AUTH}_{g+1}$ from L1.}
Then:
 \begin{itemize}
 \item Set the entry lock $Q_r$ to the highest-ranked QC
       (DATA or ANCHOR) that $r$ knows from generations
       $0,\ldots,g$.
       $Q_r$ never changes after this. QCs that arrive later do not update
       it.
 \item Enter $(g+1,0)$ \RR{as unfrozen and start its progress checks}.
 \item Send $\langle\tagg{STATUS},g+1,0,Q_r,E\rangle_r$ and its validation
       bundle to \RR{$L(g+1,0)$}. Here $E$ is the quit proof for $g$,
       which is also the entry proof of $g+1$.
 \item \RR{Process any buffered proofs for $g+1$.}
 \end{itemize}

\item \textbf{New-view.} Run this step only at $L(g+1,0)$. Wait $2\Delta$
after its own entry. If $g+1$ is still current and unfrozen:
 \begin{itemize}
 \item Let $Q^*$ be the highest-ranked QC from generations $0,\ldots,g$
       that the leader knows, including the QCs in STATUS
       messages.
 \item Let $A_{g+1}$ be the block that $Q^*$ certifies.
 \item Broadcast
       $\langle\tagg{NEWVIEW},g+1,0,A_{g+1},Q^*,E\rangle_{L(g+1,0)}$
       with its validation bundle. \RR{Choose only one opening block for
       $g+1$.}
 \end{itemize}

\item \textbf{Anchor vote.} On receiving a NEWVIEW for $g+1$:
 \begin{enumerate}[label=(\alph*)]
 \item \RR{If it contains a quit proof for $g+1$, run Quit and stop.}
 \item Validate its leader, entry proof $E$, opening block $A_{g+1}$,
       $Q^*$, and dependencies. If data are missing, request them and
       wait without voting. If the message is invalid, discard it.
 \item Vote only if all of these hold: $g+1$ is current and unfrozen; $r$
       has cast no ANCHOR vote in $g+1$; and $\rank(Q^*)\ge\rank(Q_r)$.
 \item To vote, forward the NEWVIEW to all replicas and broadcast
       $\langle\tagg{ANCHOR\mbox{-}VOTE},g+1,0,\Hash(A_{g+1})\rangle_r$.
       \RR{Start no timer. Use no fresh position.}
 \end{enumerate}
 $f+1$ matching ANCHOR votes form the ANCHOR QC. When $L(g+1,0)$
 obtains it, the leader starts Propose \RR{in \cref{alg:steady} at
 $j=0$}. Every replica, including one that cast no ANCHOR vote,
 processes DATA proposals backed by a valid ANCHOR QC under Vote.
\end{enumerate}
\end{protocol}

The entry-lock test governs ANCHOR voting only.
The delay bound $\Delta$ applies to correct transmissions.


\clearpage 


\section{Example Illustration}
\paragraph{proposing inside one term.}
\Cref{fig:steady} traces steps 1--3 of \cref{alg:steady} with $n=3$, so a
QC needs $f+1=2$ votes. The leader proposes position $j$ as soon as it
holds the QC of position $j-1$. It does not wait for that parent to commit.
Every replica starts its own $2\Delta$ timer when it votes. The timers of
several positions therefore run at the same time. Each replica commits a
block when its own timer for that block expires.

\begin{figure}[htbp]\centering
% Steady-state proposing inside one term (Protocol 1, steps 1--3).
% n = 3 and f = 1, so a QC has f+1 = 2 votes; K = 3.
% Lanes on the left, one note per event row on the right.
% Time runs downward. Message delays are drawn much shorter than Delta;
% every timer bar has the same length 2*Delta (6.2 units).
\begin{tikzpicture}[
  x=1cm,y=-1cm,font=\fontsize{9}{10.5}\selectfont,
  party/.style={draw,rounded corners=2pt,fill=white,align=center,
    inner sep=3pt,minimum width=2.4cm},
  lane/.style={draw=black!30,densely dashed},
  msg/.style={-{Stealth[length=1.7mm]},semithick},
  xchg/.style={{Stealth[length=1.3mm]}-{Stealth[length=1.3mm]},thin},
  tag/.style={fill=white,inner sep=1pt,font=\fontsize{8}{9}\selectfont},
  note/.style={anchor=north west,align=flush left,text width=6.3cm,inner sep=0pt,
    font=\fontsize{8.5}{10}\selectfont},
  dot/.style={circle,fill=black,inner sep=1.1pt},
  tA/.style={line width=2.6pt,draw=rrTeal!75},
  tB/.style={line width=2.6pt,draw=rrBlue!45}
]
  \def\xL{1.4}\def\xa{4.6}\def\xb{7.8}\def\xn{9.25}
  \path[use as bounding box] (-.3,-.55) rectangle (15.6,9.5);
  \node[party] at (\xL,0) {Leader $L(g,t)$};
  \node[party] at (\xa,0) {Replica $r_1$};
  \node[party] at (\xb,0) {Replica $r_2$};
  \foreach \x in {\xL,\xa,\xb} \draw[lane] (\x,.4)--(\x,8.45);

  % Timer bars: B_k (teal) and B_{k+1} (light blue).
  \foreach \x/\s in {\xL/1.20,\xa/1.50,\xb/1.75}
    \draw[tA] (\x+.2,\s)--(\x+.2,\s+6.2);
  \foreach \x/\s in {\xL/4.60,\xa/4.90,\xb/5.15}
    \draw[tB] (\x+.42,\s)--(\x+.42,8.45);

  % Position j = 0.
  \draw[msg] (\xL,1.20)--node[tag,pos=.5] {$B_k$} (\xa,1.50);
  \draw[msg] (\xL,1.20)--(\xb,1.75);
  \foreach \x/\y in {\xL/1.20,\xa/1.50,\xb/1.75} \node[dot] at (\x,\y) {};
  \node[note] at (\xn,.80) {\textbf{Position $j=0$.} The leader holds
    $Q_{\mathrm{parent}}$ and sends
    $\tagg{PROPOSE}(g,t,0,B_k,Q_{\mathrm{parent}})$. Each replica
    validates it, votes, and starts its own timer at $\tau_r(P)$.};

  \draw[xchg] (\xL,2.75)--(\xa,2.75);
  \draw[xchg] (\xL,2.95)--(\xb,2.95);
  \draw[xchg] (\xa,3.15)--(\xb,3.15);
  \node[note] at (\xn,2.70) {All-to-all: each replica forwards $P$ and
    multicasts $\tagg{DATA\mbox{-}VOTE}(g,t,0,\Hash(B_k))$.};

  \foreach \x in {\xL,\xa,\xb} \node[dot] at (\x,3.65) {};
  \node[note] at (\xn,3.60) {$f+1$ matching votes: every replica can form
    $\QC_{g,t}(B_k)$. A QC starts no timer.};

  % Position j = 1.
  \draw[msg] (\xL,4.60)--node[tag,pos=.5] {$B_{k+1}$} (\xa,4.90);
  \draw[msg] (\xL,4.60)--(\xb,5.15);
  \foreach \x/\y in {\xL/4.60,\xa/4.90,\xb/5.15} \node[dot] at (\x,\y) {};
  \node[note] at (\xn,4.45) {\textbf{Position $j=1$.} The leader sends
    $\tagg{PROPOSE}(g,t,1,B_{k+1},\QC_{g,t}(B_k))$ at once. The parent
    $B_k$ need not be committed. Votes follow as above.};
  \draw[xchg] (\xL,5.45)--(\xa,5.45);
  \draw[xchg] (\xL,5.60)--(\xb,5.60);
  \draw[xchg] (\xa,5.75)--(\xb,5.75);
  \foreach \x in {\xL,\xa,\xb} \node[dot] at (\x,6.0) {};

  % Position j = 2 (terminal when K = 3).
  \draw[msg] (\xL,6.25)--node[tag,pos=.5] {$B_{k+2}$} (\xa,6.55);
  \draw[msg] (\xL,6.25)--(\xb,6.80);
  \node[note] at (\xn,6.05) {\textbf{Position $j=2$.} The leader sends
    $\tagg{PROPOSE}(g,t,2,B_{k+2},\QC_{g,t}(B_{k+1}))$. With $K=3$ it is
    terminal; its QC starts rotation (\cref{fig:rotation}).};

  % B_k commits at each replica at its own vote time plus 2*Delta.
  \foreach \x/\e in {\xL/7.40,\xa/7.70,\xb/7.95} {
    \draw[thick] (\x-.12,\e)--(\x+.55,\e);
    \node[tag,anchor=east] at (\x-.12,\e) {commit $B_k$};
  }
  \node[note] at (\xn,7.85) {Each replica commits $B_k$ when its own timer
    expires, at $\tau_r(P)+2\Delta$. The timers of later positions keep
    running.};

  % Measure one timer.
  \draw[semithick] (8.6,1.75)--(8.6,7.95);
  \draw[semithick] (8.48,1.75)--(8.72,1.75);
  \draw[semithick] (8.48,7.95)--(8.72,7.95);
  \node[tag,rotate=-90] at (8.88,4.85) {$2\Delta$ timer of $B_k$ at $r_2$};

  % Legend.
  \draw[tA] (0,9.2)--(.6,9.2);
  \node[tag,anchor=west] at (.7,9.2) {timer of $B_k$};
  \draw[tB] (3.0,9.2)--(3.6,9.2);
  \node[tag,anchor=west] at (3.7,9.2) {timer of $B_{k+1}$};
  \node[tag,anchor=west] at (\xn,9.2)
    {Time runs downward. Only Quit cancels a timer.};
\end{tikzpicture}

\caption{Proposing inside term $(g,t)$ with $n=3$. Each colored bar is one
replica's local commit timer, drawn with length $2\Delta$; message delays
are drawn much shorter than $\Delta$. A QC lets the leader propose the next
position at once, so the timers of $B_k$ and $B_{k+1}$ overlap. Each
replica commits $B_k$ at its own $\tau_r(P)+2\Delta$.}
\label{fig:steady}
\end{figure}



\paragraph{Example: normal rotation.}
\Cref{fig:rotation} traces step~4 of \cref{alg:steady}. Each replica enters
$(g,t+1)$ when it first obtains the terminal QC $Q_T$, so entry times
differ. Replica $r_2$ receives votes slowly. It first sees $Q_T$ inside
the new leader's proposal. By the message-processing order, it runs
step~4 before it votes. The timers of $B_T$ keep running after the term
changes, and $B_T$ commits in term $t+1$.

\begin{figure}[htbp]\centering
% Normal rotation from term (g,t) to term (g,t+1) (Protocol 1, step 4).
% Four representative replicas. Votes sent to r_2 are slow, so r_2 first
% sees the terminal QC inside the new leader's proposal.
% Lanes on the left, one note per event row on the right.
% Time runs downward; every timer bar has length 2*Delta (6.4 units).
\begin{tikzpicture}[
  x=1cm,y=-1cm,font=\fontsize{9}{10.5}\selectfont,
  party/.style={draw,rounded corners=2pt,fill=white,align=center,
    inner sep=3pt,minimum width=2.2cm},
  lane/.style={draw=black!30,densely dashed},
  msg/.style={-{Stealth[length=1.7mm]},semithick},
  xchg/.style={{Stealth[length=1.3mm]}-{Stealth[length=1.3mm]},thin},
  tag/.style={fill=white,inner sep=1pt,font=\fontsize{8}{9}\selectfont},
  note/.style={anchor=north west,align=flush left,text width=6.0cm,inner sep=0pt,
    font=\fontsize{8.5}{10}\selectfont},
  dot/.style={circle,fill=black,inner sep=1.1pt},
  tA/.style={line width=2.6pt,draw=rrTeal!75},
  tB/.style={line width=2.6pt,draw=rrBlue!45},
  entry/.style={draw=rrBlue,line width=1.2pt,dotted}
]
  \def\xO{1.2}\def\xa{3.7}\def\xb{6.2}\def\xN{8.7}\def\xn{9.55}
  \path[use as bounding box] (-.4,-.7) rectangle (15.6,9.45);
  \node[party] at (\xO,0) {Old leader\\$L(g,t)$};
  \node[party] at (\xa,0) {Replica\\$r_1$};
  \node[party] at (\xb,0) {Replica\\$r_2$};
  \node[party] at (\xN,0) {New leader\\$L(g,t+1)$};
  \foreach \x in {\xO,\xa,\xb,\xN} \draw[lane] (\x,.55)--(\x,8.75);

  % Local entry to (g,t+1): each replica enters when it first obtains Q_T.
  \draw[entry] (-.2,3.6)--(\xO,3.6)--(\xa,3.6)--(\xb,5.5)--(\xN,3.6)--(9.3,3.6);

  % Timers: B_T keeps running across the term boundary.
  \foreach \x/\s in {\xO/1.20,\xa/1.40,\xb/1.60,\xN/1.80}
    \draw[tA] (\x+.2,\s)--(\x+.2,\s+6.4);
  \foreach \x/\s in {\xN/5.30,\xb/5.50,\xa/5.70,\xO/5.90}
    \draw[tB] (\x+.42,\s)--(\x+.42,8.75);

  % Terminal proposal of term (g,t).
  \draw[msg] (\xO,1.20)--node[tag,pos=.5] {$B_T$} (\xa,1.40);
  \draw[msg] (\xO,1.20)--(\xb,1.60);
  \draw[msg] (\xO,1.20)--(\xN,1.80);
  \foreach \x/\y in {\xO/1.20,\xa/1.40,\xb/1.60,\xN/1.80} \node[dot] at (\x,\y) {};
  \node[note] at (\xn,1.05) {The old leader sends
    $\tagg{PROPOSE}(g,t,K{-}1,B_T,\ldots)$ for the terminal position.
    Every replica votes and starts a timer for $B_T$.};

  \draw[xchg] (\xO,2.75)--(\xa,2.75);
  \draw[xchg] (\xa,2.95)--(\xN,2.95);
  \draw[xchg] (\xO,3.15)--(\xN,3.15);
  \node[note] at (\xn,2.70) {All-to-all votes. The votes sent \emph{to}
    $r_2$ happen to be slow.};

  % Obtain Q_T and run step 4.
  \foreach \x in {\xO,\xa,\xN} \node[dot] at (\x,3.6) {};
  \node[note] at (\xn,3.55) {Three replicas obtain the terminal QC
    $Q_T=\QC_{g,t}(B_T)$. Step~4: each enters $(g,t{+}1)$, records its
    entry time, and forwards $Q_T$ to the new leader.};
  \draw[msg] (\xa,4.05)--node[tag,pos=.78] {backup $Q_T$} (\xN,4.30);

  % The new leader's first proposal carries Q_T.
  \draw[msg] (\xN,5.30)--node[tag,pos=.5] {$B_{T+1}$} (\xb,5.50);
  \draw[msg] (\xN,5.30)--(\xa,5.70);
  \draw[msg] (\xN,5.30)--(\xO,5.90);
  \foreach \x/\y in {\xN/5.30,\xb/5.50,\xa/5.70,\xO/5.90} \node[dot] at (\x,\y) {};
  \node[note] at (\xn,5.15) {The new leader sends
    $\tagg{PROPOSE}(g,t{+}1,0,B_{T+1},Q_T)$, with $Q_T$ as its parent QC.};
  \node[note] at (\xn,6.30) {$r_2$ first sees $Q_T$ inside this proposal.
    It runs step~4 first and enters $(g,t{+}1)$. Then it votes for
    $B_{T+1}$.};

  % B_T commits at every replica, after its local term boundary.
  \foreach \x/\e in {\xO/7.60,\xa/7.80,\xb/8.00,\xN/8.20} {
    \draw[thick] (\x-.12,\e)--(\x+.55,\e);
    \node[tag,anchor=east] at (\x-.12,\e) {commit $B_T$};
  }
  \node[note] at (\xn,7.75) {$B_T$ commits at every replica in term
    $t{+}1$. Rotation cancels no timer.};

  % Legend.
  \draw[tA] (-.1,9.15)--(.5,9.15);
  \node[tag,anchor=west] at (.6,9.15) {timer of $B_T$};
  \draw[tB] (2.6,9.15)--(3.2,9.15);
  \node[tag,anchor=west] at (3.3,9.15) {timer of $B_{T+1}$};
  \draw[entry] (5.3,9.15)--(5.9,9.15);
  \node[tag,anchor=west] at (6.0,9.15) {local entry to $(g,t{+}1)$};
  \node[tag,anchor=west] at (\xn,9.15) {Time runs downward.};
\end{tikzpicture}

\caption{Normal rotation from $(g,t)$ to $(g,t+1)$. The dotted line marks
each replica's local entry to $(g,t+1)$. Every holder of $Q_T$ forwards it
to the new leader as a backup. The new leader's first proposal carries
$Q_T$ as its parent QC. Rotation cancels no timer: $B_T$ commits at every
replica after that replica has entered $(g,t+1)$.}
\label{fig:rotation}
\end{figure}




\paragraph{Example: quit from different terms.}
\Cref{fig:quit} traces \cref{alg:recovery} with $n=3$, $f=1$, and $K=1$.
Correct $H_0$ leads $(g,t)$, Byzantine $F$ leads $(g,t+1)$, and correct
$H_1$ leads $(g+1,0)$. $H_0$ obtains the terminal QC $Q_T$ quickly and
enters $(g,t+1)$. $H_1$ is still in $(g,t)$. $F$ then signs two different
blocks for $(g,t+1,0)$, which is a type (b) quit proof. Both correct
replicas freeze all of $g$, even though they are in different terms. Quit
cancels the timers of $B_T$, so $B_T$ is not yet committed. Recovery keeps
$B_T$: its QC $Q_T$ becomes every entry lock, the new leader re-certifies
$B_T$ with an ANCHOR QC, and the first fresh child commits $B_T$ as an
ancestor.

\begin{figure}[htbp]\centering
% Quit and recovery when correct replicas freeze from different terms
% (Protocol 2). n = 3 and f = 1, so a QC or quit proof needs 2 signatures.
% K = 1. H_0 leads (g,t), Byzantine F leads (g,t+1), H_1 leads (g+1,0).
% Lanes on the left, one note per event row on the right.
% Time runs downward; vertical spacing shows causal order, not scale.
\begin{tikzpicture}[
  x=1cm,y=-1cm,font=\fontsize{9}{10.5}\selectfont,
  party/.style={draw,rounded corners=2pt,fill=white,align=center,
    inner sep=3pt,minimum width=2.3cm},
  lane/.style={draw=black!30,densely dashed},
  msg/.style={-{Stealth[length=1.7mm]},semithick},
  rd/.style={-{Stealth[length=1.7mm]},semithick,densely dashed},
  xchg/.style={{Stealth[length=1.3mm]}-{Stealth[length=1.3mm]},semithick},
  tag/.style={fill=white,inner sep=1pt,font=\fontsize{8}{9}\selectfont},
  note/.style={anchor=north west,align=flush left,text width=5.55cm,inner sep=0pt,
    font=\fontsize{8.5}{10}\selectfont},
  phase/.style={anchor=west,fill=white,inner sep=2pt,
    font=\fontsize{9}{10.5}\selectfont\bfseries},
  band/.style={draw,rounded corners=2pt,fill=black!4,align=center,
    inner sep=3pt,font=\fontsize{8.5}{10}\selectfont},
  dot/.style={circle,fill=black,inner sep=1.1pt},
  tA/.style={line width=2.6pt,draw=rrTeal!75},
  wt/.style={line width=2.6pt,draw=rrBlue!70!black,dash pattern=on 2.5pt off 1.5pt},
  cut/.style={text=rrRed,inner sep=0pt,font=\bfseries\fontsize{12}{12}\selectfont}
]
  \def\xA{1.2}\def\xF{3.9}\def\xB{6.6}\def\xC{9.1}\def\xn{9.95}
  \path[use as bounding box] (-.4,-.7) rectangle (15.6,20.45);
  \node[party] at (\xA,0) {$H_0$ (correct)\\$=L(g,t)$};
  \node[party,fill=black!8] at (\xF,0) {$F$ (Byzantine)\\$=L(g,t+1)$};
  \node[party] at (\xB,0) {$H_1$ (correct)\\$=L(g+1,0)$};
  \node[party,minimum width=1.5cm] at (\xC,0) {L1};
  \foreach \x in {\xA,\xF,\xB,\xC} \draw[lane] (\x,.55)--(\x,20.05);

  % Local clocks: commit timers (teal), cancelled timers (red cross),
  % and waits (dashed blue).
  \draw[tA] (\xA+.2,1.30)--(\xA+.2,6.65);
  \node[cut] at (\xA+.2,6.65) {$\times$};
  \draw[tA] (\xB+.2,3.75)--(\xB+.2,8.30);
  \node[cut] at (\xB+.2,8.30) {$\times$};
  \draw[wt] (\xA+.2,12.15)--(\xA+.2,13.20);
  \draw[wt] (\xB+.2,12.25)--(\xB+.2,13.10);
  \draw[wt] (\xB+.2,15.50)--(\xB+.2,16.50);
  \draw[tA] (\xB+.2,18.90)--(\xB+.2,19.90);
  \draw[tA] (\xA+.2,19.05)--(\xA+.2,19.90);
  \foreach \x in {\xA,\xB} \draw[thick] (\x-.12,19.90)--(\x+.55,19.90);

  % ---------------------------------------------------------------
  \node[phase] at (-.35,.85) {1.\ Term $(g,t)$ ends at different times for $H_0$ and $H_1$};
  \draw[msg] (\xA,1.30)--node[tag,pos=.5] {$B_T$} (\xF,1.45);
  \draw[msg] (\xA,1.30)--node[tag,pos=.78] {$B_T$ (slow)} (\xB,3.75);
  \node[dot] at (\xA,1.30) {};
  \node[note] at (\xn,1.20) {$H_0$ proposes $B_T$ at $(g,t,0)$. With
    $K=1$, it is the terminal block. The copy to $H_1$ is slow.};
  \draw[msg] (\xF,1.80)--node[tag,pos=.5] {vote} (\xA,1.95);
  \node[dot] at (\xA,1.95) {};
  \node[tag,anchor=east] at (\xA-.1,1.95) {$Q_T$};
  \node[note] at (\xn,2.35) {$F$ votes at once. $H_0$ has $f+1=2$ votes:
    the terminal QC $Q_T$. Step~4: $H_0$ enters $(g,t{+}1)$.};
  \node[dot] at (\xB,3.75) {};
  \node[note] at (\xn,3.65) {$H_1$ receives $B_T$, votes, and starts its
    timer. It has no $Q_T$, so it stays in $(g,t)$.};

  % ---------------------------------------------------------------
  \node[phase] at (-.35,5.0) {2.\ $F$ equivocates; both correct replicas freeze all of $g$};
  \draw[msg] (\xF,5.45)--node[tag,pos=.5] {$B$} (\xA,5.60);
  \draw[msg] (\xF,5.75)--node[tag,pos=.5] {$B'$} (\xA,5.90);
  \node[note] at (\xn,5.40) {$F$ sends two different signed blocks
    $B\neq B'$ for $(g,t{+}1,0)$. Together they form a type (b) quit
    proof.};
  \node[dot] at (\xA,6.65) {};
  \node[note] at (\xn,6.60) {Evidence first. $H_0$ runs Quit: it freezes
    $g$, cancels the timer of $B_T$, records $a_{H_0}$, and broadcasts the
    proof.};
  \draw[msg] (\xA,8.10)--node[tag,pos=.75] {quit proof} (\xB,8.30);
  \node[dot] at (\xB,8.30) {};
  \node[note] at (\xn,8.10) {$H_1$ is still in $(g,t)$, but the proof is
    valid. $H_1$ also runs Quit and records $a_{H_1}$. It keeps $Q_T$ from
    the proof's admission witness.};
  \draw[msg] (\xB,9.75)--node[tag,pos=.5] {proof} (\xC,9.90);
  \node[note] at (\xn,9.70) {A correct replica submits the proof to L1.
    L1 later excludes $F$.};
  \node[band,text width=9.3cm] at (4.55,10.95)
    {At the freeze, $H_0$ is in term $(g,t{+}1)$ and $H_1$ is still in
     $(g,t)$.\\ $B_T$ is certified by $Q_T$ but not yet committed.};

  % ---------------------------------------------------------------
  \node[phase] at (-.35,11.75) {3.\ Recovery into $(g+1,0)$};
  \node[note] at (\xn,12.10) {Each waits until
    $\max\{a_r{+}\Delta,\,d(g{+}1)\}$, then reads $\tagg{AUTH}_{g+1}$
    from L1.};
  \draw[rd] (\xC,13.00)--(\xB,13.10);
  \node[tag] at (7.95,12.72) {read $\tagg{AUTH}_{g+1}$};
  \draw[rd] (\xC,13.00)--(\xA,13.20);
  \foreach \x/\y in {\xA/13.55,\xB/13.55} \node[dot] at (\x,\y) {};
  \node[note] at (\xn,13.45) {Each sets its entry lock $Q_r=Q_T$, the
    highest QC it knows, and enters $(g{+}1,0)$.};
  \draw[msg] (\xA,14.70)--node[tag,pos=.75] {STATUS} (\xB,14.85);
  \node[note] at (\xn,14.60) {$H_0$ sends $\tagg{STATUS}(g{+}1,0,Q_T,E)$ to
    the new leader $H_1$.};
  \node[note] at (\xn,15.50) {$H_1$ waits $2\Delta$ from its own entry.
    $Q^*=Q_T$ is the highest QC it knows.};
  \draw[msg] (\xB,16.55)--(\xF,16.65);
  \draw[msg] (\xB,16.55)--node[tag,pos=.75] {NEWVIEW} (\xA,16.75);
  \node[note] at (\xn,16.45) {$H_1$ sends
    $\tagg{NEWVIEW}(g{+}1,0,A_{g+1},Q_T,E)$, with opening block
    $A_{g+1}=B_T$.};
  \draw[xchg] (\xA,17.75)--node[tag,pos=.5] {ANCHOR votes} (\xB,17.75);
  \node[note] at (\xn,17.60) {$\rank(Q^*)\ge\rank(Q_r)$ holds, so both send
    ANCHOR votes. $f+1$ votes form the ANCHOR QC. It starts no timer.};
  \draw[msg] (\xB,18.90)--node[tag,pos=.5] {$B_{T+1}$} (\xA,19.05);
  \node[dot] at (\xB,18.90) {}; \node[dot] at (\xA,19.05) {};
  \node[note] at (\xn,18.85) {$H_1$ proposes the first fresh child
    $B_{T+1}$. Both vote and start timers. On expiry, each commits
    $B_{T+1}$ and its ancestor $B_T$.};
\end{tikzpicture}

\caption{Quit and recovery when correct replicas freeze from different
terms. Teal bars are commit timers; a red cross marks a timer that Quit
cancels. Dashed blue bars are waits. The ANCHOR QC starts no timer; the first fresh block $B_{T+1}$
starts the timers that commit $B_T$. Vertical spacing shows causal order,
not delay to scale.}
\label{fig:quit}
\end{figure}

\section{Execution times and progress deadlines}
\label{sec:timing}
\begin{lemma}[Generation entry and freeze propagation]
\label{lem:entry}
Correct entry times to $(g,0)$ differ by at most $\Delta$. If the first
correct freeze of $g$ is at $a$, every correct replica freezes it by
$a+\Delta$. All recover to $(g+1,0)$, with entry skew at most $\Delta$.
\end{lemma}
\begin{proof}
The initial skew is assumed. A proof is broadcast on first freeze.
Any correct replica already in $g$ receives it by $a+\Delta$.
A replica still finishing prior recovery enters $g$ within $\Delta$
of its first correct entry, hence by $a+\Delta$, and processes the
buffered proof. No correct replica can complete recovery from $g$ before
the first correct freeze of $g$.

All replicas use the same public availability time $d(g+1)$. Hence
\[
\begin{aligned}
 |e_r(g+1,0)-e_s(g+1,0)|
 &=\big|\max\{a_r+\Delta,d(g+1)\}
          -\max\{a_s+\Delta,d(g+1)\}\big|\\
 &\le |a_r-a_s|\le\Delta.
\end{aligned}
\]
This gives the next generation's entry bound directly.
A buffered proof may freeze the next generation immediately after entry;
it does not erase that entry event or change its target.
\end{proof}

\begin{lemma}[Terminal-QC handoff]
\label{lem:handoff}
Let $q$ be the first time a correct replica in generation $g$ has verified
a complete terminal QC for $(g,t)$. If the successor is correct and the
generation stays unfrozen, it proposes its first child by $q+\Delta$,
and all correct replicas process that proposal by $q+2\Delta$.
Without assuming a correct successor, all correct replicas still obtain
the terminal QC by $q+2\Delta$ if the generation remains unfrozen.
\end{lemma}
\begin{proof}
The first correct QC holder sends the complete bundle to the successor.
It arrives by $q+\Delta$. If that leader is still finishing generation
entry, \cref{lem:entry} places its entry no later than $q+\Delta$.
It has the signed terminal block and its QC and immediately proposes
the child, unless it already did so. Its proposal reaches all by
$q+2\Delta$.

For the final claim, a correct terminal-QC signer voted at some $s\le q$
and broadcast its complete proposal. All correct replicas have entered
the generation and process it by $s+\Delta$. Earlier-term proposals
remain eligible, so all cast or have already cast that DATA vote.
Their $f+1$ or more correct votes give the terminal QC everywhere by
$s+2\Delta\le q+2\Delta$.
\end{proof}
Here $q$ is a correct replica's verified QC receipt, not an adversary's
private assembly time. Ordinary term-entry skew is at most $2\Delta$
in this unfrozen execution, whereas recovery-generation entry skew is
$\Delta$. A common clock does not require simultaneous term entry.

\begin{figure}[htbp]\centering
\begin{tikzpicture}[x=1cm,y=-1cm,font=\small,
 msg/.style={-{Stealth[length=2mm]},semithick,draw=rrBlue},
 lab/.style={fill=white,inner sep=3pt,align=center}]
\path[use as bounding box] (0,-.5) rectangle (14.8,5.0);
\node[align=center] at (1.9,0) {First correct holder\\of terminal QC};
\node[align=center] at (7.35,0) {Correct next leader\\$L(g,t+1)$};
\node[align=center] at (12.6,0) {Other correct\\replicas};
\foreach \x in {1.9,7.35,12.6} \draw[black!30,dashed] (\x,.6)--(\x,4.4);
\node[anchor=east] at (1.7,1.15) {$q$};
\draw[msg] (1.9,1.15)--node[lab,above,sloped]{terminal QC + complete history}(7.35,2.05);
\node[lab,anchor=west] at (7.5,2.05) {by $q+\Delta$: enter; propose};
\draw[msg] (7.35,2.65)--node[lab,above,sloped]{first DATA proposal + parent QC}(12.6,3.65);
\draw[msg] (7.35,2.65)--(1.9,3.65);
\node[lab] at (7.35,4.35) {By $q+2\Delta$: all have processed the new proposal.};
\end{tikzpicture}
\caption{Normal handoff while the generation remains unfrozen. The arrows
show causal stages, not mandatory waiting. Each replica records its own
first terminal-QC entry time. Backup forwarding introduces no vote phase.}
\label{fig:handoff}
\end{figure}

\begin{lemma}[Correct-leader progress]
\label{lem:progress}
While a correct leader's term remains active and no correct replica has
frozen the generation, its $i$th fresh position triggers every correct
replica's DATA vote by that replica's deadline
$e_r(g,t)+(2i+4)\Delta$, for $i=1,\ldots,K$.
Thus a correct replica sends no progress BLAME against such a leader
before the first correct freeze.
\end{lemma}
\begin{proof}
For recovered term zero, put $e=\min_r e_r(g,0)$. The correct leader enters
by $e+\Delta$ and finishes its wait by $e+3\Delta$. All correct STATUS
messages arrive within that wait, so its selected QC meets every entry
lock. NEWVIEW reaches all by $e+4\Delta$; ANCHOR votes reach the leader
by $e+5\Delta$; its first fresh proposal reaches all by $e+6\Delta$.
Each further proposal costs at most $2\Delta$: votes back to the leader,
then the next proposal to all. The $i$th vote therefore occurs by
\[
 e+6\Delta+2(i-1)\Delta=e+(2i+4)\Delta
 \le e_r(g,0)+(2i+4)\Delta.
\]
Genesis has no ANCHOR exchange and is faster.

For $t>0$, let $q$ be as in \cref{lem:handoff} for the preceding terminal
QC. The first proposal arrives by $q+2\Delta$ and the $i$th by
$q+2i\Delta$. Every local entry is at or after $q$, so these arrivals
precede the stated deadlines. History proposals are processed even after
ordinary advancement, and completed-term checks are disabled. At equality,
received proposals are processed before a progress callback.
\end{proof}

\begin{figure}[htbp]\centering
\begin{tikzpicture}[node distance=4mm,font=\small,
 event/.style={draw=rrBlue!65,rounded corners=2pt,align=left,
 text width=12.2cm,inner sep=5pt},
 arr/.style={-{Stealth[length=1.7mm]},draw=rrBlue}]
\node[event] (a) {$e$: earliest correct entry to recovered term zero};
\node[event,below=of a] (b) {$e+\Delta$: latest possible correct leader entry};
\node[event,below=of b] (c) {$e+3\Delta$: leader completes its own $2\Delta$ wait; sends NEWVIEW};
\node[event,below=of c] (d) {$e+4\Delta$: all correct replicas send ANCHOR votes};
\node[event,below=of d] (f) {$e+5\Delta$: leader has ANCHOR QC; sends first fresh DATA proposal};
\node[event,below=of f,draw=rrTeal] (h) {$e+6\Delta$: all correct replicas vote; first progress check is met};
\foreach \u/\v in {a/b,b/c,c/d,d/f,f/h} \draw[arr] (\u.south)--(\v.north);
\end{tikzpicture}
\caption{Worst-case recovery timeline for a correct initializer. Subsequent
fresh positions add $2\Delta$ each, giving $(2i+4)\Delta$. ANCHOR votes
are not fresh progress. Uninterrupted commit timers commit the first
block at all correct replicas by $e+8\Delta$.}
\label{fig:recovery}
\end{figure}
\FloatBarrier

\section{Safety analysis}
\label{sec:safety}
\begin{lemma}[One generation, one chain or a proof]
\label{lem:fork}
Two admitted histories in the same generation have the same opening and
are prefix-comparable, or their signed material gives a valid quit proof.
\end{lemma}
\begin{proof}
Different openings give two conflicting signed NEWVIEWs. With one opening,
the height formula assigns each fresh height to exactly one $(t,j)$ and
leader. The first divergent fresh height then gives two different signed
DATA proposals at that position. The histories provide their admission
witnesses. An older-generation fork in an inherited justification is not
by itself a new-generation equivocation.
\end{proof}

\begin{lemma}[A completed commit timer protects its block]
\label{lem:timer}
If a correct replica votes for $B$ at $s$ and commits it at $s+2\Delta$,
every correct replica knows its DATA QC by that expiry. We call this
QC the \emph{protecting QC} of $B$. No conflicting
QC in the same generation, nor an ANCHOR QC for a different opening,
can exist, including a QC assembled later from old votes.
\end{lemma}
\begin{proof}
A correct freeze by $s+\Delta$ would reach the committer by $s+2\Delta$
and cancel its timer. Thus none occurs that early. The complete proposal
sent at $s$ reaches every correct replica by $s+\Delta$; even a replica
still finishing recovery has entered this generation by then, by
\cref{lem:entry}. Any incompatible retained history would expose a proof
and freeze that receiver by $s+\Delta$, a contradiction. Therefore every
correct replica has already cast or now casts the same DATA vote. Normal
term advancement and prior indirect commitment do not suppress that vote.
Their at least $f+1$ multicast votes provide the QC everywhere by expiry.

Any conflicting QC has a correct signer. If it processed $B$ first, its
retained evidence prevents the conflicting vote. Otherwise its conflicting
vote and forwarded signed history preceded that processing, no later
than $s+\Delta$. That history reaches the committer by expiry and exposes
the fork, contradicting commitment. An ANCHOR vote for another opening
similarly forwards the conflicting NEWVIEW. This reasoning bounds the
correct vote's time, not the later time of QC assembly. Evidence-first
handling covers simultaneous arrival and expiry.
\end{proof}

\begin{lemma}[Recovery locks preserve a completed timer]
\label{lem:lock}
If a correct timer commits $B$ in generation $g$, every correct entry
lock in a later generation ranks at least as high as its protecting
QC.
\end{lemma}
\begin{proof}
Let $s$ be its vote time and $a$ the first correct freeze of $g$.
The preceding cancellation argument gives $a>s+\Delta$. Every successor
entry satisfies
$e_r(g+1,0)\ge a_r+\Delta\ge a+\Delta>s+2\Delta$.
The protecting QC is therefore known before every lock snapshot.
Highest-QC selection and retained locks preserve the rank bound in all
later generations. Additional L1 waiting only delays these snapshots.
If the generation never freezes, there is no later entry to check.
\end{proof}

\begin{lemma}[Unique extensibility]
\label{lem:extension}
For a block $B$ committed by a correct timer, every admitted QC ranked
at least as high as its protecting QC certifies an extension of $B$.
\end{lemma}
\begin{proof}
In its generation, \cref{lem:timer} excludes alternative certified
openings and conflicting QCs. With a fixed opening, higher terms own
strictly higher fresh heights, so a sufficiently ranked QC cannot certify
a strict ancestor of $B$.

If a later-generation counterexample exists, choose one of least rank.
An ANCHOR QC has a correct voter whose older justification meets its
entry lock, hence the protecting rank by \cref{lem:lock}. That justification
certifies the same nonextending block at a strictly smaller rank, a
contradiction. A DATA QC has a lower-ranked parent QC in the same later
generation, including at a term boundary. Its parent cannot extend $B$
if the child does not, so this too is a smaller counterexample. The
lexicographic nonnegative ranks are well-ordered; content tie-breaks do
not affect these strict descents.
\end{proof}

\begin{theorem}[Safety and validity]
\label{thm:safety}
Under \cref{ass:network,ass:authority}, correct committed logs are
prefix-comparable, never revoked, and application-valid.
\end{theorem}
\begin{proof}
Call a commitment \emph{direct} when a block's own timer makes it, and
\emph{indirect} when it commits an ancestor. Order the protecting QCs of
any two direct  commitments by rank.
\Cref{lem:extension} makes the higher-ranked block extend the lower-ranked
one. Indirect ancestor commitments inherit comparability. Commitments
are never withdrawn by either transition. Complete validation establishes
application validity. Neither proposer exclusion nor schedule fairness
is used in this argument.
\end{proof}

\section{Liveness analysis}
\label{sec:liveness}
\begin{lemma}[One timer rule suffices]
\label{lem:window}
Suppose a correct replica votes for and forwards a DATA proposal at $s$.
If no correct replica freezes that generation through $s+3\Delta$, every
correct replica commits its block by $s+3\Delta$; the originating timer
commits at $s+2\Delta$.
\end{lemma}
\begin{proof}
Every correct replica has entered and processes the proposal by $s+\Delta$.
An incompatible history would force a freeze, contrary to the premise.
Thus every correct replica votes by that time or has already voted.
Every such vote starts the one commit timer. Earlier terms and
QC-carrying history use exactly the same rule. All timers expire by
$s+3\Delta$ without cancellation; prior commitments already satisfy
the conclusion. No QC-start timer is needed.
\end{proof}

\begin{lemma}[Only finitely many generations expose equivocation]
\label{lem:finite}
Under \cref{ass:exclusion}, only finitely many generation indices are
named by admissible equivocation proofs revealed to correct replicas.
\end{lemma}
\begin{proof}
Each identified leader is Byzantine, since correct leaders sign one
opening and one DATA block per position. For each identity $p$ ever exposed,
its first revealed proof names $g_{\mathrm{bad}}$ and accountability gives
a finite $G_p>g_{\mathrm{bad}}$. Any admissible equivocation by $p$ must
belong to a generation $g<G_p$: at $g\ge G_p$, the public authority
never assigns $p$ a leader position. The union of these finite ranges
over the finite exposed identity set is finite. A proof revealed later
still names its original generation and cannot restart a completed recovery.
\end{proof}
This does not bound interruptions by $f$. An identity can cause several
before exclusion becomes effective.

\begin{lemma}[A correct initializer obtains a commit window]
\label{lem:initializer}
Consider a generation with correct $L(g,0)$ in which no admissible
equivocation proof is revealed. Its first fresh proposal and local DATA
vote occur at a finite time $s$, and every correct replica commits that block by
$s+3\Delta$, even if QC rotation reaches a Byzantine successor.
\end{lemma}
\begin{proof}
Let $e$ be earliest entry. Before the first fresh proposal, no later
term is admitted: its terminal-QC ancestry would require that proposal.
The earliest possible correct term-zero BLAME is at $e+6\Delta$.
With equivocation excluded, no quit proof can freeze this initial interval.
By the timeline of \cref{lem:progress} (\cref{fig:recovery}), the
initializer sends its first fresh proposal by $e+5\Delta$.
 Thus $s$ exists. Genesis initialization is faster.

Suppose the first correct freeze were at $a\le s+3\Delta$.
It must use BLAME. Before that first freeze, \cref{lem:progress} excludes
any correct term-zero BLAME. For a later term, a correct BLAME signer
entered at some $u\ge s$: its admission history includes the
initializer's first proposal, signed at $s$. The first complaint is no
earlier than $u+6\Delta\ge s+6\Delta$, even if Byzantine BLAME
signatures were prepared in advance. Thus no complete BLAME proof can
cause the supposed freeze. \Cref{lem:window} completes the claim.
\end{proof}

\begin{theorem}[Eventual committed-height progress]
\label{thm:height}
Under \cref{ass:network,ass:authority,ass:exclusion,ass:fairness}, every
correct replica eventually commits arbitrarily large heights.
\end{theorem}
\begin{proof}
If infinitely many generations are entered, \cref{lem:finite} leaves a
suffix without revealed equivocation. Schedule fairness supplies infinitely
many correct initializers there. By \cref{lem:initializer}, each commits
its first fresh block everywhere. Its certified opening extends prior
committed blocks by locks and unique extensibility, so this fresh child
has strictly greater height. Authority availability prevents a permanent
wait for any one next generation.

If only finitely many generations are entered, the last one cannot freeze.
After a freeze, every correct replica enters the next generation at
the finite time $\max\{a_r+\Delta,d(g+1)\}$ (\cref{ass:authority}), which
would be another entry.
Suppose this unfrozen generation had a last term. Every observed
terminal QC spreads through its correct signer's proposal and the
resulting correct votes, by \cref{lem:handoff}. Thus all correct replicas
eventually reach that last common term. If any correct replica votes for
its terminal proposal, forwarding and correct votes produce the terminal
QC everywhere and admit another term. Otherwise every correct replica has
fewer than $K$ votes at its last progress check and broadcasts BLAME.
Their $f+1$ or more matching BLAME messages form a quit proof and
force a freeze. Both contradict a last term.

Hence certified terms and heights increase without bound. Each terminal
QC has a correct signer who forwards its complete proposal; all correct
replicas vote under the same rule and their timers commit, since this
generation never freezes. This proves progress at every correct replica.
\end{proof}

\begin{theorem}[Request liveness]
\label{thm:requests}
Under the same assumptions, every persistently admissible correct-client
request is eventually committed by all correct replicas.
\end{theorem}
\begin{proof}
In the infinite-generation case, infinitely many correct initializers
after the last equivocation generation commit their first fresh blocks
everywhere. In the finite-generation case, term fairness gives infinitely
many correct leaders in the last unfrozen generation; their fresh blocks
commit. Fix a request after dissemination. Only finitely many requests
precede it. Each such successful correct proposal includes the oldest
fitting pending requests or commits them as ancestors. Speculative
inclusion does not remove a pending request. Induction over the finite
predecessors therefore commits the chosen request, and safety preserves it.
\end{proof}

\section{Attacks on tempting simplifications and weaker assumptions}
\label{sec:attacks}
The preceding theorems concern the specified protocol and assumptions.
The following executions identify exactly what fails when a rule or
assumption is removed. They are included to distinguish a proved conditional
guarantee from an unconditional claim.

\paragraph{A. Skipping earlier-term proposals without QC-start
timers.}
Take $n=3$, $f=1$, $K=1$, correct replicas $H_0,H_1$, and Byzantine $F$.
In each generation let $H_0$ initialize and $F$ follow it; later terms
can cycle fairly over all identities. $H_0$ proposes $B$ and starts its
timer. $F$ supplies its vote, forms a terminal QC, and gives that QC to
$H_1$ before $H_1$ processes the proposal. If $H_1$ first advances and
then refuses earlier-term proposals, it never starts $B$'s timer.
$F$ stays silent. $H_0$ commits. Timeout recovery later
re-certifies $B$ with an ANCHOR QC, which starts no commit
timer.
Repeating this execution can leave $H_1$ without any non-genesis
commitment. There is no equivocation to exclude.
\Cref{alg:steady} prevents this execution by processing unseen history
proposals and preserving earlier-term voting eligibility.

\paragraph{B. Shorten normal-term deadlines to $2i\Delta$.}
Use the same committee and $K=1$. Let $H_0$ propose at $s$ and self-vote.
Choose $0<\varepsilon<\Delta/3$. Fast exchange with $F$ gives $H_0$ the
terminal QC and successor entry at $u=s+2\varepsilon$, while $H_1$
first receives the proposal and votes at $s+\Delta$. $F$, the successor,
sends no DATA proposal and gives its signed BLAME to $H_0$ in
advance. $H_0$ commits at $s+2\Delta$, then under the shortened rule sends
BLAME and freezes at $u+2\Delta$. Deliver the resulting proof to $H_1$
within $\varepsilon$. It arrives at
\[
 s+2\Delta+3\varepsilon<s+3\Delta,
\]
canceling $H_1$'s timer before its expiry. All delays obey the upper bound.
Recovery re-certifies the block with an ANCHOR QC and starts no
timer. Repeat with $H_0$ as initializer:
$H_0$ commits increasing heights but $H_1$ never commits a non-genesis
block. $F$ only withholds, so exclusion does not stop the attack.
Thus a deadline sufficient to avoid blaming a correct successor need not
preserve all-replica liveness. The specified $6\Delta$ first check leaves
the $3\Delta$ window of \cref{lem:initializer} intact.

\paragraph{C. Remove permanent exclusion.}
For any fixed finite $K$, let a correct initializer and $F$ exchange votes
with a sufficiently small positive delay $\varepsilon$. They can certify
all $K$ blocks and admit $F$ as successor before the first $2\Delta$
timer expires. $F$ signs two different valid first children of the terminal
block. Deliver their equivocation proof to all correct replicas before
any timer expires. They cancel every timer and recover. Repeat in every
generation with the same still-eligible $F$. QCs and heights can grow
while no non-genesis block commits. A finite term length $K$,
including $K=n$, does not exclude this execution when the model supplies
no lower bound that forces a long enough proposal window. After the first
revelation identifies $g_{\mathrm{bad}}$, the specified finite $G_p$
prohibits reusing the same leader in every later generation; the
execution removes that condition.

\paragraph{D. Start the handoff bound at private QC formation.}
Let $F$ lead a $K=1$ term, with correct successor $H_0$ and another correct
replica $H_1$, and choose $0<\varepsilon<\Delta$.
$H_0$ receives the terminal proposal and votes at $s$.
$F$ privately combines that vote with its own at $q=s+\varepsilon$,
but withholds its vote and QC. Deliver $H_0$'s forwarded proposal and vote
to $H_1$ at $s+\Delta$, and $H_1$'s vote and backup QC to $H_0$ at
$s+2\Delta$.
$H_0$ cannot obtain the terminal QC before $s+2\Delta>q+\Delta$.
The bound in \cref{lem:handoff} instead starts at a correct replica's
complete verified QC receipt, which triggers actual forwarding.

\paragraph{E. Treat two unauthenticated messages as equivocation.}
An outsider can fabricate different blocks for the same position and
force arbitrary freezes if signatures and admission are not checked.
Evidence-first processing therefore prioritizes verification of a proof;
it does not bypass that verification. Conversely, a malformed enclosing
message must not suppress a valid proof already present in its signed
components.

\section{Communication and latency}
\label{sec:complexity}
We count only complete, admissible messages sent by correct replicas.
We exclude Byzantine flooding and retrieval traffic, discussed
below.
Each correct replica forwards and votes for at most one DATA
block per position, sends at most one BLAME  per term, and
broadcasts its first quit proof once per generation. Let $b$ bound a fresh
payload, $\sigma$ a signature, and $D$ a bound on additional
dependency bytes in each normalized bundle charged to the operation.
Different operations may have different $D$. Explicit QCs have
$O(n\sigma)$ bytes.

\begin{center}\small
\begin{tabularx}{\linewidth}{@{}>{\raggedright\arraybackslash}p{.24\linewidth}
 >{\raggedright\arraybackslash}p{.20\linewidth}Y@{}}
\toprule
Operation & Correct messages & Conservative bytes\\
\midrule
DATA forwarding and votes & $O(n^2)$ per position &
$O(n^2(b+D+n\sigma))$\\
Terminal-QC backup & $O(n)$ per certified terminal &
$O(n(b+D+n\sigma))$ to one successor\\
Progress complaints & $O(n^2)$ per admitted term &
$O(n^2\sigma)$ plus admission/context data\\
Quit and recovery & $O(n^2)$ per recovery &
$O(n^2(D+n\sigma))$ for proof, STATUS, NEWVIEW, and ANCHOR votes\\
Local commit timers & No messages & No extra vote or QC phase\\
L1 bulletin board & External service & Evidence submission, public authority
reads, and eligibility updates\\
\bottomrule
\end{tabularx}
\end{center}
Different correct replicas can initially accept different
equivocations, but each still sends only one DATA vote/forward per position
before proof processing freezes the generation. Certificate subsets do not
defeat semantic deduplication. Let $S$ count distinct processed DATA
positions $(g,t,j)$, $T$ distinct admitted terms $(g,t)$ including
unfinished ones, and $R$ generations whose recovery has been initiated,
including those still waiting for L1. The message bound is
\[
 O\bigl(n^2(S+T+R)+nT\bigr).
\]
Partial BLAME traffic is charged even when no recovery follows. For
normally completed $K$-block terms, backup forwarding costs $O(n/K)$
messages per block; complaints cost at most $O(n^2/K)$ per block.
Votes and old DATA processed during recovery remain charged to their
originating positions. The bulletin-board model requires no BFT relay of
L1 finality proofs; its service costs are separate from these message counts.

\paragraph{History and retrieval costs.}
The displayed bounds cover complete bundles. Malformed or incomplete
Byzantine traffic can induce arbitrary requests; it is outside this
bounded-traffic accounting. Any implementation that sends cache references
must separately count requests, replies, and transferred missing objects,
and preserve the complete-message timing assumption. $D$ is not constant:
without cache reuse, forwarding length-$\ell$ prefixes can require
$D=O(\ell(b+n\sigma))$ plus control evidence. Across $S$ successive
positions, repeated full-prefix transmission can cost
$O(n^2S^2(b+n\sigma))$ bytes, even though DATA message counts remain
$O(n^2S)$. No linear amortized history-byte claim or constant-size QC
assumption is made.

\paragraph{Normal operation.}
A proposal accepted at $\tau_r(P)$ commits locally at
$\tau_r(P)+2\Delta$ if its timer survives. From a correct vote at $s$,
an uninterrupted window gives commitment everywhere by $s+3\Delta$.
From the first correct terminal-QC time $q$, a correct successor proposes
by $q+\Delta$, all receive by $q+2\Delta$, and surviving timers commit
that child everywhere by $q+4\Delta$. With homogeneous actual delay
$\delta$, ordinary proposal/QC iterations take about $2\delta$.
The backup adds no mandatory $\Delta$ wait: a leader acts as soon as
its parent QC is available. It does add the stated $O(n)$ terminal traffic.

\Needspace{13\baselineskip}
\paragraph{Recovery.}
Let $a=\min_r a_r$ be the first correct quit time. The common L1
availability time gives
\[
 e=\min_r e_r(g+1,0)=\max\{a+\Delta,d(g+1)\}.
\]
With a correct initializer, the first DATA proposal reaches all by
$e+6\Delta$ and an uninterrupted window commits it everywhere by
\[
 \max\{a+\Delta,d(g+1)\}+8\Delta.
\]
If $d(g+1)\le a+\Delta$, this is $a+9\Delta$; otherwise recovery
waits for the public authority. Eventual exclusion and
authority availability give no numerical bound on the last equivocation,
number of aborted generations, or end-to-end recovery delay. Quantitative
bounds require corresponding L1 and scheduling bounds.

\begin{thebibliography}{9}
\bibitem{sync}
I. Abraham, D. Malkhi, K. Nayak, L. Ren, and M. Yin.
\emph{Sync HotStuff: Simple and Practical Synchronous State Machine Replication}.
IEEE Symposium on Security and Privacy, 2020, pp. 106--118.
\href{https://conferences.computer.org/sp/pdfs/sp/2020/349700a106.pdf}{Published paper}.
\bibitem{schneider}
F. B. Schneider. Implementing fault-tolerant services using the state machine
approach: a tutorial. \emph{ACM Computing Surveys}, 22(4):299--319, 1990.
\href{https://doi.org/10.1145/98163.98167}{DOI}.
\end{thebibliography}
\end{document}
